% !TeX root = ../../progress-sem1.tex
% 1.x ขอบเขตของโครงงาน
\section{ขอบเขตของโครงงาน}

\thaijustify{
  \begin{enumerate}
    \item ใช้โมเดล Generative AI ที่มีอยู่ในปัจจุบันผ่าน CPE AI Gateway เพื่อสร้าง two-alternative UI prototype จากการอธิบายความต้องการทางการออกแบบด้วยภาษาธรรมชาติ
    \item จัดเตรียมสภาพแวดล้อมสำหรับให้ผู้เข้าร่วมทดสอบ comparative usability testing ตรวจสอบ เปรียบเทียบ และให้ข้อเสนอแนะเชิงบริบท
    \item ใช้โมเดล AI ที่มีอยู่ในปัจจุบันผ่าน CPE AI Gateway ในการรวบรวมข้อเสนอแนะและระบุรูปแบบ ข้อมูลเชิงลึก และประเด็นสำคัญเพื่อพิจารณา โดยให้นักออกแบบเป็นผู้ตัดสินใจ (human-in-the-loop) เมื่อแนวโน้มความคิดเห็นของผู้เข้าร่วมทดสอบมีความผิดปกติ
    \item ใช้การวิเคราะห์ที่ได้จาก AI ในการเชื่อมโยงข้อเสนอแนะของผู้ใช้เข้ากับการเปลี่ยนแปลงของการออกแบบและสร้างรายงาน
    \item แพลตฟอร์มนี้ใช้ภาษาอังกฤษเป็นหลัก เนื่องจากข้อจำกัดของการประมวลผลภาษาไทยของ Generative AI
  \end{enumerate}
}

\subsection{คุณสมบัติของซอฟต์แวร์ (Software Features)}
\thaijustify{
  แพลตฟอร์มแบ่งคุณสมบัติออกเป็นกลุ่มตามหลักการ Design Thinking ซึ่งประกอบด้วยขั้น Empathize, Define, Ideate, Prototype และ Test~\cite{trantruong2025genaiux} โดยมีส่วนพื้นฐานที่ใช้ร่วมกันทุกขั้น ดังนี้
  \begin{enumerate}
    \item \textbf{ส่วนพื้นฐาน (Core)}
      \begin{itemize}
        \item \textbf{การยืนยันตัวตนและสิทธิ์ตามบทบาท}\\
          สมัครสมาชิกและเข้าสู่ระบบด้วยบัญชี Google หรืออีเมลและรหัสผ่าน ยืนยันอีเมล รีเซ็ตรหัสผ่าน และกำหนดสิทธิ์การเข้าถึงตามบทบาทในแต่ละโปรเจกต์
        \item \textbf{การตั้งค่าโปรเจกต์}\\
          สร้าง workspace และเชิญสมาชิก สร้างโปรเจกต์และกำหนดบทบาทของสมาชิก
      \end{itemize}
    \item \textbf{Empathize, Define และ Ideate}
      \begin{itemize}
        \item \textbf{ความต้องการทางการออกแบบ}\\
          กำหนดความต้องการทางการออกแบบ (design intent) ได้แก่ persona, design system, มาตรฐานการเข้าถึง, สีแบรนด์, แนวทาง CI และ tone of voice เพื่อใช้เป็น context ของ AI ในทุกรอบ
      \end{itemize}
    \item \textbf{Prototype}
      \begin{itemize}
        \item \textbf{การสร้าง UI ด้วย AI}\\
          เขียน user story, problem statement และเลือกอุปกรณ์ที่ต้องการทดสอบ ตรวจและแก้ structured prompt ที่ AI ร่าง สร้าง two-alternative UI (variant A และ B) พร้อมติดตามขั้นตอนการทำงานของ AI แบบ real-time ปรับแก้ UI ผ่านแชตหรือเลือก element บนหน้าจอ และย้อนดูเวอร์ชันก่อนหน้า
      \end{itemize}
    \item \textbf{Test}
      \begin{itemize}
        \item \textbf{พื้นที่ทดสอบ (Testing space)}\\
          กำหนดภารกิจสำหรับผู้เข้าร่วมทดสอบ คำถามเพื่อเปรียบเทียบ สมมติฐานพร้อมคำถามเฉพาะจุด และจำนวนรอบสูงสุดต่อผู้เข้าร่วมทดสอบ ผู้เข้าร่วมทดสอบเข้าเซสชันด้วยรหัสหรือลิงก์โดยไม่ต้องมีบัญชี ให้ความยินยอมตาม PDPA ทำภารกิจบนทั้งสองดีไซน์ ตอบคำถามผ่านแชต และปักหมุดความคิดเห็นบน component หรือให้ความเห็นต่อทั้งหน้าจอ
      \end{itemize}
    \item \textbf{Empathize ของวงรอบถัดไป}
      \begin{itemize}
        \item \textbf{การวิเคราะห์ข้อเสนอแนะ}\\
          ติดตามความคืบหน้าและแนวโน้มจำนวนความคิดเห็นของผู้เข้าร่วมทดสอบแต่ละคน AI จัดกลุ่มข้อเสนอแนะเป็นประเด็น (theme) ให้นักออกแบบตรวจ ปรับป้ายกำกับ รวม แยก หรือระบุว่าไม่แก้ไขพร้อมเหตุผล แล้วนำประเด็นที่เลือกไปสร้าง UI สำหรับรอบถัดไป รวมถึงสร้าง Persona Artifact จากข้อมูลการทดสอบ
        \item \textbf{รายงานสำหรับผู้มีส่วนได้ส่วนเสีย}\\
          สร้างรายงานแยกตามบทบาท ได้แก่ รายงานผลสำหรับ Client, เอกสารส่งมอบงาน (handoff) สำหรับ Developer และความคืบหน้าของรอบการทดสอบสำหรับ Project Manager โดยทุกการเปลี่ยนแปลงในรายงานสืบย้อนกลับไปยังข้อเสนอแนะต้นทางได้ ผู้มีส่วนได้ส่วนเสียรับทราบและแสดงความเห็นต่อรายงานได้ และส่งออกเป็น PDF หรือลิงก์
      \end{itemize}
  \end{enumerate}
}

\subsection{กลุ่มเป้าหมาย}
\thaijustify{
  กลุ่มเป้าหมายหลักของแพลตฟอร์มคือนักออกแบบ UX/UI ที่ทำงานเป็นทีมหรือในสตูดิโอออกแบบ และต้องการทดสอบการใช้งานกับผู้ใช้จริงบ่อยขึ้นภายใต้กรอบเวลาส่งมอบงานที่จำกัด ส่วนกลุ่มเป้าหมายรองคือผู้มีส่วนได้ส่วนเสียในการพัฒนาซอฟต์แวร์ ผู้ใช้แบ่งตามบทบาทได้ดังนี้
  \begin{itemize}
    \item \textbf{Designer}\\
      ผู้ใช้หลัก กำหนดความต้องการทางการออกแบบ สร้างและปรับแก้ UI ด้วย AI จัดเซสชันการทดสอบ ตรวจผลการวิเคราะห์ของ AI และตัดสินใจขั้นสุดท้ายว่าจะปรับปรุงการออกแบบอย่างไร
    \item \textbf{Usability Tester}\\
      ผู้เข้าร่วมทดสอบที่นักออกแบบเชิญมา เป็นตัวแทนของผู้ใช้ปลายทางของผลิตภัณฑ์ ไม่จำเป็นต้องมีบัญชีในแพลตฟอร์ม เพียงกรอกชื่อและอีเมลเมื่อเข้าเซสชันด้วยรหัสหรือลิงก์
    \item \textbf{Client}\\
      ผู้ว่าจ้างหรือเจ้าของผลิตภัณฑ์ ดูรอบการออกแบบที่อนุมัติแล้วและรายงานผลการทดสอบ รับทราบและแสดงความเห็นต่อรายงาน
    \item \textbf{Developer}\\
      นักพัฒนาที่รับช่วงงานออกแบบไปพัฒนาต่อ ดูเอกสารส่งมอบงาน ได้แก่ spec, component, โค้ด และความเสี่ยงทางเทคนิค
    \item \textbf{Project Manager}\\
      ผู้จัดการโครงการ ติดตามความคืบหน้าของรอบการทดสอบ กำหนดการ และทรัพยากรที่ใช้
  \end{itemize}
}

\subsection{สภาพแวดล้อม}
\thaijustify{
  แพลตฟอร์มพัฒนาเป็นเว็บแอปพลิเคชันที่ใช้งานผ่านเว็บเบราว์เซอร์ โดยไม่ต้องติดตั้งโปรแกรมเพิ่มเติม ผู้ใช้ทุกบทบาทใช้งานผ่านเบราว์เซอร์บนคอมพิวเตอร์ ส่วนผู้เข้าร่วมทดสอบเข้าสู่เซสชันการทดสอบด้วยรหัสหรือลิงก์ สภาพแวดล้อมของระบบมีรายละเอียดดังนี้
  \begin{itemize}
    \item \textbf{ฝั่งผู้ใช้}\\
      เว็บเบราว์เซอร์รุ่นปัจจุบัน ส่วนหน้าจอที่ AI สร้างขึ้นรองรับการแสดงผลสามแบบ ได้แก่ desktop, tablet และ mobile ตามที่กำหนดในแต่ละโปรเจกต์
    \item \textbf{ฝั่งเซิร์ฟเวอร์}\\
      Node.js รุ่น 20.19 ขึ้นไป ใช้ฐานข้อมูล PostgreSQL 17 และ Redis 7 ซึ่งรันผ่าน Docker
    \item \textbf{บริการ AI}\\
      เรียกใช้โมเดลภาษาขนาดใหญ่ผ่าน CPE AI Gateway ของภาควิชาวิศวกรรมคอมพิวเตอร์ มหาวิทยาลัยเทคโนโลยีพระจอมเกล้าธนบุรี ซึ่งเข้าถึงได้ผ่านเครือข่าย VPN ของมหาวิทยาลัยเท่านั้น
    \item \textbf{การนำระบบขึ้นใช้งาน}\\
      \todo{ระบุเซิร์ฟเวอร์หรือบริการที่ใช้ deploy ระบบ (ใน repository ยังไม่มีการตั้งค่า deploy)}
  \end{itemize}
}
